\section{Implications for the Search for Technosignatures}\label{technosignatures-implications}

\subsection{Reframing the Drake Equation's Longevity Term}\label{reframing-drake}

The Drake Equation \citep{Drake1965-wp} estimates the number of detectable civilizations in our galaxy as
\begin{equation}
N = R_*\, f_p\, n_e\, f_l\, f_i\, f_c\, L ,
\label{eq:drake}
\end{equation}
where the final term \emph{L}, the average duration that civilizations remain detectable, carries outsized importance and remains deeply uncertain. Our simulations suggest that the traditional interpretation of $L$ as a continuous span of emission is unrealistic. To account for the intermittent nature of civilizational activity, we introduce an effective detectability duration,
\begin{equation}
L_{\mathrm{eff}} = D_c \, T_{\mathrm{span}} ,
\label{eq:leff}
\end{equation}
where \(D_{c}\) is the mean duty cycle and \(T_{\mathrm{span}}\) is the observation window (1000 years in our simulations). If $L$ is better understood as $L_{\mathrm{eff}}$, then models that assume continuous detectability may significantly overestimate $N$. A galaxy populated by civilizations with low duty cycles would remain largely silent, even if intelligent life is common. We note that a related but distinct concept, a ``reappearance factor'' for planets where civilizations can arise multiple times in succession \citep{Bhattacharya2011-qr}, which counts sequential independent civilizations over geological timescales, whereas our duty cycle refers to the active vs.\ inactive phases within a single civilization's lifespan.

Moreover, technology can outlast its creators, as the longevity of a technosignature may far exceed the lifespan of the civilization that produced it \citep{Wright2022-wh,Cirkovic2019-oq}. This further complicates the interpretation of $L$, which must account not only for active emission but also for the persistence of relic technosignatures, as we examine below.

These considerations help explain the Great Silence. Civilizations may exist in abundance, but if they are detectable only intermittently, the probability that their technosignatures overlap with our observational window is low \citep{Balbi2018-fp}. For two intermittent civilizations, the probability of mutual visibility is the product of their duty cycles, which can be very small. The simulation results raise a deeper question: should we expect civilizations to be uniform in character, or drawn from a diverse ensemble of trajectories? If most civilizations follow fragile, collapse-prone paths (e.g., S1, S2, S6), then long-duty-cycle civilizations would be statistical outliers, and the apparent absence of signals may reflect not the rarity of extraterrestrial technology but the dominance of intermittently silent technospheres. Our findings support this view: only S3 and S10 avoided collapse entirely, while most scenarios exhibited high collapse frequency and moderate or low duty cycles. These collapse-prone trajectories suggest that civilizations may not follow monotonic energy-growth trajectories, but instead oscillate or regress, challenging the static assumptions of the Kardashev scale \citep{Kardashev1964-gv}.

Recent work by \citet{Balbi2024-bg} provides further theoretical support for this view. Applying Lindy's law to technosignature longevities, they show that for energy-intensive technoemissions the longevity distribution follows a power law $\rho_L(L) \propto L^{-\alpha}$ with $\alpha > 1$, strongly disfavoring long-lived technosignatures. Under realistic assumptions ($\alpha \geq 2$), the first detected technosignature is predicted to be relatively short-lived ($L < 10^2$ years at 95\% confidence). If most detectable technosignatures are intrinsically short-lived, then the combination of Lindy-distributed longevities and low duty cycles (as observed in our collapse-prone scenarios) implies a doubly diminished probability of detection, offering a complementary explanation for the Great Silence.

\subsection{Technosignature Persistence and Observation Probabilities}\label{technosignature-persistence}

Eq. (\ref{eq:leff}) gives the detectability duration in terms of the mean duty cycle and time in the observation window. Some technosignatures may have a lifetime exactly equal to the the civilization's lifetime, but others may persist for some time after collapse. For technosignature $i$, the detectability duration is given as
\begin{equation}
    L_{i}=D_cT_{\text{span}} + \text{min}[T_i, (1-D_c)T_{\text{span}}]\delta_{i},
\end{equation}
where $T_i$ is the lifetime of technosignature $i$ and $\delta_{i}$ is a flag that is 1 if technosignature $i$ is present in the scenario and 0 otherwise.

The probability that technosignature $i$ could be observed (assuming an appropriate observatory) is
\begin{equation}
    p_{i}=\frac{L_i}{T_{\text{span}}} = D_c + \text{min}[\frac{T_i}{T_{\text{span}}}, 1-D_c]\delta_{i},
\end{equation}
These probabilities are shown in the figure above for four example technosignatures, with a 1000-year time span. These probabilities should be interpreted as the probability of finding a given technosignature for remote observations of Earth in a given scenario, assuming the use of a capable observatory and sufficient observation time. Nitrogen dioxide (NO$_2$) has a short lifetime of hours to days, so the probability is about equal to the duty cycle for scenarios in which nitrogen dioxide is present. The lifetimes of CFC-11 is 55\,yr, and the lifetime of CFC-12 is 140\,yr, which could be observed in five of the scenarios. CFC-12 has higher observation probabilities due to its longer lifetime. Carbon tetrafluoride (CF$_4$) has a lifetime of 1000\,yr or longer, observable in one of the scenarios.

%\begin{figure*}[t]
%\centering
%\includegraphics[width=\textwidth]{../figures/technosignature_dectability.png}
%\caption{he probability that a technosignature could be observed, taking into account the duty cycle for each scenario and the lifetime of each technosignature.}
%\label{fig:ts_detect}
%\end{figure*}

The $L_{\mathrm{eff}}$ formulation assumes that detectability ends when activity ceases. As noted above, this need not be the case \citep{Balbi2021-bc}. The persistence of different technosignature types varies enormously. Orbital debris and artificial satellites can last from decades to as long as 10,000 to 100,000 years. Satellites in low Earth orbit typically deorbit within decades due to atmospheric drag, while those in geosynchronous orbit, facing minimal resistance, may persist for much longer. Without maintenance, a dense band of satellites in the so-called Clarke belt would slowly degrade through collisions and orbital perturbations over tens of thousands of years, eventually becoming undetectable \citep{Socas-Navarro2018-kg}. Megastructures such as Dyson spheres, space mirrors, or stellar engines may persist longer still, though their survival depends on structural design and orbital stability. Rigid constructs without active station-keeping may collapse or drift within a few centuries to several thousand years \citep{Wright2020-be}, however, more modular configurations, like swarms of solar collectors, could remain in stable orbits for far longer timescales \citep{Smith2022-ds}.

Atmospheric technosignatures provide the shortest windows of detectability. On Earth, CFC-11 and CFC-12 have atmospheric lifetimes of about 55 and 140 years, respectively, before being broken down by photolysis or other chemical processes \citep{NOAA-GML-lz}. Other compounds are more resilient: tetrafluoromethane (CF$_4$) persists for \textasciitilde50,000 years, hexafluoroethane (C$_2$F$_6$) for \textasciitilde10,000 years, and octafluoropropane (C$_3$F$_8$) for \textasciitilde2,600 years \citep{Trudinger2016-ky}. Longer-lived isotopic signatures are conceivable, but their remote detection remains technically challenging. Electromagnetic technosignatures such as radio or laser transmissions propagate indefinitely but are only practically detectable for finite periods. Focused, high-power signals might remain observable for a few hundred light-years over spans of centuries to millennia, depending on the strength of the signal and the sensitivity of receivers \citep{Sheikh2025-td}.

These persistence times imply that technosignature lifetimes can range from negligible (radio bursts) to dominant (orbital infrastructure). For short-lived civilizations with brief ON phases, persistent technosignatures could substantially extend the effective detectability window beyond what $L_{\mathrm{eff}}$ alone would suggest. In contrast, for continuously active civilizations, persistence contributes relatively little. If most detectable technosignatures are short-lived (e.g., radio leakage), then the duty cycle remains the key parameter, and our results suggest low duty cycles make detection unlikely. But if many technosignatures persist for millennia beyond collapse, then the silence we observe becomes harder to attribute solely to civilizational fragility; it would require that either technosignatures are intrinsically ephemeral, or civilizations rarely produce persistent artifacts.

In summary, our results suggest that if detectability is often sporadic, then long periods without contact may be the expected baseline. Intermittency may be typical, and recognizing that could refine our expectations and expand the scope of SETI. We also note that some persistent technologies, particularly autonomous systems such as Bracewell probes \citep{Bracewell1960-bp} or orbital maintenance infrastructure, may not decay passively but instead cycle through their own active and dormant states, creating layered duty-cycle structures that merit further investigation.
